\section{Problem Formulation}
\label{sec:problem}

\noindent\textbf{Tokenizers and token cost.}
Let $\mathcal{L}$ be a set of $L$ languages. A BPE vocabulary $V$ consists of a base
alphabet and an ordered list of merges, with size bounded by a budget $K$. Each language
$\ell$ has a training text $\mathcal{D}_\ell$ from which merges are learned, and
$c_\ell(a,b)$ denotes the number of occurrences of the adjacent pair $(a,b)$ in
$\mathcal{D}_\ell$ under the current vocabulary. A multi-parallel development set contains
the same $S$ sentences in every language. Let $T_\ell(V)$ be the number of tokens that $V$ produces on the development
sentences of language $\ell$. We define the \emph{token cost} of language $\ell$ as
\begin{equation}
  t_\ell(V) \;=\; \frac{T_\ell(V)}{S},
  \label{eq:cost}
\end{equation}
the average number of tokens needed for one sentence of shared content. A tokenizer is fair
when all $t_\ell$ are close to one another, and efficient when they are small.

\noindent\textbf{Choice of cost unit.}
The denominator in Eq.~\eqref{eq:cost} must not depend on the script: in UTF-8 a Devanagari
character occupies three bytes and a Latin character one, so per byte English appears to be
the most expensive language even when it needs the fewest tokens per sentence, and any rule
that helps the worst-off language then sends its merges to English
(Section~\ref{sec:unitresult}). We therefore measure cost per parallel sentence, following
\citet{petrov2023} and \citet{foroutan2026}.

\noindent\textbf{A multi-objective view.}
The costs $\mathbf{t}(V)=(t_1(V),\dots,t_L(V))$ form a vector of objectives to be made small
together:
\begin{equation}
  \min_{|V|\le K}\; \mathbf{t}(V).
  \label{eq:mo}
\end{equation}
A vocabulary is Pareto-optimal for Eq.~\eqref{eq:mo} if no other vocabulary of size at most
$K$ lowers the cost of one language without raising that of another, and the cost vectors
of Pareto-optimal vocabularies form the compression--parity front. Existing tokenizers each
commit to one point on it: classical BPE favours total compression, Parity-aware BPE the
worst-off language.

\noindent\textbf{Power-mean scalarisation.}
To trace the front with a single parameter, we scalarise Eq.~\eqref{eq:mo} with a weighted
power mean,
\begin{equation}
  \min_{|V|\le K}\; F_p(V) \;=\;
  \Bigl(\sum_{\ell\in\mathcal{L}} \pi_\ell\, t_\ell(V)^{p}\Bigr)^{1/p},
  \qquad p\in[1,\infty],
  \label{eq:obj}
\end{equation}
with weights $\pi_\ell>0$ that sum to one. We set $\pi_\ell$ to the share of language $\ell$
in the training text, measured in development-sentence units, for the reason given in
Section~\ref{sec:rule}. The exponent $p$ sets how strongly expensive
languages are emphasised: at $p=1$, $F_1=\sum_\ell \pi_\ell t_\ell$ is the weighted average
cost that classical BPE reduces; as $p\to\infty$, $F_p\to\max_\ell t_\ell$, the cost of the
worst-off language that Parity-aware BPE reduces; for $1<p<\infty$, $F_p$ lies between the
two and gives more weight to a language the more expensive it is. For finite $p$ and positive weights, $F_p$ is strictly increasing in every $t_\ell$.
Consequently, an exact minimiser of Eq.~\eqref{eq:obj} is Pareto-optimal for
Eq.~\eqref{eq:mo}, while a minimiser of the max at $p=\infty$ is only guaranteed to be
weakly Pareto-optimal. In practice, we minimise Eq.~\eqref{eq:obj} greedily, one merge at a
time (Section~\ref{sec:method}), so the curve we measure is the front of the greedy
algorithm rather than the exact Pareto front.

\noindent\textbf{Evaluation measures.}
For parity, we report the Gini coefficient of the per-language costs on a held-out parallel
test set,
\begin{equation}
  G \;=\; \frac{\sum_{i,j\in\mathcal{L}} |t_i - t_j|}{2L^2\,\bar t},
  \label{eq:gini}
\end{equation}
where $\bar t$ is the mean cost ($G=0$ when every language costs the same). For efficiency,
we report compression in bytes per token on the training text, the quantity classical BPE
greedily maximises, and additionally on the balanced test set.